\documentclass[11pt]{article}
\usepackage[margin=1in]{geometry}
\usepackage[T1]{fontenc}
\usepackage[utf8]{inputenc}
\usepackage{lmodern}
\usepackage{microtype}
\usepackage{amsmath}
\usepackage{booktabs}
\usepackage{hyperref}
\title{Automatically Aggregated Drum2Grid Results}
\author{}
\date{}
\begin{document}
\maketitle
\section{Methodology}
\subsection{Dataset and codec configuration}
The diffusion cache contains 11523 training, 1534 validation, and 1733 test windows. Audio is represented with DAC at 44100 Hz using 9 quantizers. The target token rate is 86.1328 Hz, and the primary target is a 72-dimensional framewise PCA projection of the full 1024-dimensional summed DAC RVQ codebook-embedding trajectory. The modeled target is the post-quantizer summed embedding, not the raw pre-quantizer encoder output.

\subsection{Seconds-aware symbolic front-end}
The conditioning front-end is the hybrid seconds-aware local front-end with radii \texttt{[0, 22, 41, 55]}, 64-dimensional per-scale features, and \texttt{seconds} positional encoding. It samples symbolic drum-grid context at DAC token times, so symbolic conditioning and codec-latent targets remain aligned under BPM-derived timing.
\begin{center}\small\begin{tabular}{ll}\toprule Component & Value \\ \midrule Symbolic grid rate & 250 Hz \\ Drum families & 8 \\ Numeric channels & 24 \\ Articulation encoding & per-family IDs to one-hot lanes \\ Radii & 0, 22, 41, 55 grid steps \\ Radii in ms & 0, 88, 164, 220 ms \\ Branch output & 64 \\ Concatenated conditioning dim & 256 \\ Denoiser width & 768 \\ \bottomrule\end{tabular}\end{center}
Radius 0 denotes the feature sample at the codec-frame time itself; nonzero radii add symmetric local windows around that time.

\subsection{PCA-latent diffusion model}
The completed diffusion runs model 72-dimensional normalized PCA coefficients with a 768-wide Transformer, 6 layers, and 8 attention heads. Plain PCA diffusion runs are available for denoising step counts 6, 12, 25, 50; auxiliary RVQ cross-entropy runs are available for step counts 6, 12, 25. Each stored quantitative row uses one generated sample per conditioning input; the paired confidence intervals therefore summarize clip variation, not sampling-seed variation.

\subsection{Compared systems}
The aggregation includes reconstruction ceilings, non-learned symbolic/rendering baselines, a direct PCA regressor, plain PCA diffusion runs, and PCA diffusion runs with auxiliary RVQ cross-entropy when their artifacts are available. Rows with missing metric artifacts are retained and flagged in \texttt{missing\_metrics.csv} rather than silently dropped.
The symbolic grid renderer is a deterministic procedural renderer from the eight-family grid, with event times from the 250 Hz seconds grid and velocities from the onset/state lanes. The symbolic nearest-neighbor baseline retrieves from the training split using normalized dot products over flattened symbolic features plus BPM, then decodes the retrieved train clip's DAC source codes.
The direct PCA regressor uses the same conditioning frontend and PCA target family but a deterministic Transformer head. Its width is larger than the diffusion model, so the comparison is conservative with respect to capacity but not architecture-identical.

\subsection{Statistical protocol}
Paired clip-level intervals use 2,000 percentile bootstrap resamples at 95\% confidence. Best-versus-rest paired tests use 2,000 two-sided sign-flip permutations over shared dataset indices, with Holm correction within each metric. FAD uncertainty is reported from repeated extrapolation runs and is not used as a paired significance test.

\subsection{Reproducible result aggregation}
All reported tables are generated by \texttt{scripts/build\_paper\_results.py}. The script gathers completed run artifacts without rerunning inference, writes normalized CSV/JSON files, and emits the standalone \LaTeX{} tables used in this section. The aggregated artifact contains 13 systems.

\section{Results and Analysis}
Table~\ref{tab:auto_main_results} is generated directly from the normalized run manifest. It should be interpreted as an automatic-metric summary: distributional quality, paired spectral/onset fidelity, direct reconstruction, and runtime are separate axes.

\begin{table*}[t]
\centering
\small
\setlength{\tabcolsep}{3pt}
\begin{tabular}{llrrrrrrrr}
\toprule
System & Type & Clips & $\mathrm{FAD}_{\infty}$ $\downarrow$ & $\mathrm{FAD}\text{-}R^2$ $\uparrow$ & Mel $\downarrow$ & Flux $\uparrow$ & MRSTFT $\downarrow$ & Audio $L_1$ $\downarrow$ & RTF $\downarrow$ \\
\midrule
Target DAC reconstruction & reconstruction\_ceiling & 1733 & 0.016 & 0.848 & 0.10 & 0.999 & 0.001 & 0.0002 & 0.015 \\
Target PCA reconstruction & reconstruction\_ceiling & 1733 & 0.016 & 0.847 & 0.10 & 0.999 & 0.001 & 0.0002 & 0.015 \\
Symbolic grid render & baseline & 1733 & 0.551 & 0.147 & 19.27 & 0.763 & 0.264 & 0.0859 & 0.004 \\
Source-code decode & baseline & 1733 & 0.016 & 0.848 & 0.10 & 0.999 & 0.001 & 0.0002 & 0.015 \\
Symbolic NN retrieval & baseline & 1733 & 0.025 & 0.846 & 17.54 & 0.330 & 0.208 & 0.0671 & 0.018 \\
Direct PCA regressor & direct\_pca & 1733 & 0.355 & 0.169 & 13.04 & 0.836 & 0.136 & 0.0451 & 0.021 \\
Diffusion PCA 6 steps & diffusion\_pca & 1733 & 0.023 & 0.851 & 6.39 & 0.843 & 0.108 & 0.0534 & 0.060 \\
Diffusion PCA 12 steps & diffusion\_pca & 1733 & 0.021 & 0.835 & 5.75 & 0.850 & 0.106 & 0.0533 & 0.065 \\
Diffusion PCA 25 steps & diffusion\_pca & 1733 & 0.019 & 0.891 & 5.69 & 0.848 & 0.106 & 0.0536 & 0.077 \\
Diffusion PCA 50 steps & diffusion\_pca & 1733 & 0.024 & 0.744 & 5.71 & 0.839 & 0.109 & 0.0544 & 0.107 \\
Diffusion PCA+RVQ-CE 6 steps & diffusion\_pca\_rvq\_ce & 1733 & 0.022 & 0.830 & 5.47 & 0.866 & 0.105 & 0.0509 & 0.027 \\
Diffusion PCA+RVQ-CE 12 steps & diffusion\_pca\_rvq\_ce & 1733 & 0.020 & 0.845 & 5.39 & 0.864 & 0.106 & 0.0515 & 0.046 \\
Diffusion PCA+RVQ-CE 25 steps & diffusion\_pca\_rvq\_ce & 1733 & 0.020 & 0.848 & 5.47 & 0.863 & 0.104 & 0.0512 & 0.072 \\
\bottomrule
\end{tabular}
\caption{Automatically aggregated test-set metrics. Missing cells indicate that the corresponding evaluation artifact was not present when the table was generated. $\mathrm{FAD}_{\infty}$ is distributional; $\mathrm{FAD}\text{-}R^2$ is the extrapolation-fit diagnostic; the paired metrics quantify reconstruction and symbolic-control-related structure.}
\label{tab:auto_main_results}
\end{table*}

\paragraph{Main observations.}
The current artifacts support four controlled observations.
First, Diffusion PCA+RVQ-CE 12 steps gives the lowest learned-system mel error (5.39 dB), a 58.6\% relative reduction from the direct PCA regressor (13.04 dB).
Second, Diffusion PCA+RVQ-CE 6 steps gives the strongest onset-flux agreement among learned systems (0.866), supporting the auxiliary RVQ-CE term as codec-structure supervision rather than a discrete-token generator.
Third, Diffusion PCA+RVQ-CE 25 steps gives the best learned-system MRSTFT (0.104), a 23.6\% relative reduction from direct PCA regression.
Fourth, phase-sensitive waveform $L_1$ is the main counterexample: Direct PCA regressor has the lowest learned-system waveform $L_1$ (0.0451), so the diffusion results should be interpreted as spectral/transient improvements rather than universal samplewise reconstruction dominance.
All completed diffusion rows are faster than real time in the stored run; RTF differences should be interpreted as implementation/runtime measurements rather than architectural consequences of RVQ-CE.
The best completed diffusion $\mathrm{FAD}_{\infty}$ row is Diffusion PCA 25 steps (0.019).

\end{document}
